\documentclass{article}
\usepackage{graphicx} % Required for inserting images
\usepackage{stil}
\usepackage{booktabs} % for nicer tables
\usepackage{array}    % for column customization
\usepackage{longtable} % optional, for multipage tables
\usepackage{caption}  % to customize captions, if needed

\title{Trade, Migration, and Global Challenges}
\author{Econ VM A}
\date{winter term 2025/2026}

\graphicspath{{syllabus/}}

\begin{document}

\maketitle
\thispagestyle{empty}

\section*{General Information}
\begin{tabular}{@{}ll}
Responsible unit:  & Economics Department \\
Instructor:          & Freya Rubel \\
Office hours: & by appointment (drop me an email)\\
Contact:           & f.rubel@fu-berlin.de \\
Class time:    & Thursday, 10:00am-12:00pm, room: 201 \\
Credits:          & 7 ECTS for active participation and examination \\
& 3 ECTS for active participation\\
Language:       & English \\
\end{tabular}

\section*{Course Description}
The course aims to introduce students to the interplay between trade, migration, and issues such as inequality, environmental degradation, and global politics. It should equip students with the analytical tools and essential knowledge necessary to engage in contemporary debates on these issues, helping them to gain a deeper understanding of the interconnected dynamics in the globalized economy. The classes will cover both theoretical concepts and empirical studies. By the end of the course, students are expected to:
\begin{itemize}
    \item Understand the key theories and concepts underlying international trade and migration.
    \item Have developed an awareness of how trade and migration trends are influenced by and contribute to challenges in areas like the labour market, environmental sustainability, or social cohesion.
\end{itemize}
For the final examination, students will be expected to demonstrate their understanding of theoretical concepts and models discussed throughout the course, as well as their ability to critically discuss the relationship between trade or migration and at least one of the challenges covered in course.
\section*{Requirements}
Students are expected to have passed the Econ OM (orientation module), which consists of the courses \textit{Economic Reasoning} and \textit{Statistical Reasoning} or be able to demonstrate that they have acquired other equivalent economics and statistics skills at other institutions or through other study programs.  \\
Active participation during the class (attendance of at least 80\%) is mandatory. A student presentation is required for active participation as well.
\section*{Evaluation}
For those who choose to be examined in this course, the final grade will be made up of the in-class student presentation (25\%) and an exam (75\%) at the end of the semester. \\
The list of papers to present will be published as soon as the final number of participating students has been determined. 
\section*{Course Outline\footnote{Can be subject to changes. Please check your emails and the FU blackboard regularly.}}
\thispagestyle{empty}

\begin{center}
\begin{tabular}{>{\bfseries}c l}
\toprule
\textbf{Session} & \textbf{Date} \\
\midrule
\midrule
\multicolumn{2}{l}{\textbf{\Roman{1} Organisation and Introduction to Globalization}} \\
\midrule
1 & October 16  \\ 
2 & October 23 \\ 
\midrule
\multicolumn{2}{l}{\textbf{\Roman{2} Trade}} \\
\midrule
3 & October 30 \\ % Freya
4 & November 6 \\ % Louise 
5 & November 13 \\ % Louise 
6 & November 20 \\ % Freya
7 & November 27 \\ % Freya 
\midrule
\multicolumn{2}{l}{\textbf{\Roman{3} Migration}} \\
\midrule
8 & December 4 \\  % Louise 
9 & December 11 \\ % Freya 
10 & December 18 \\ % Louise
11 & January 8 \\ % Freya 
\midrule
\multicolumn{2}{l}{\textbf{\Roman{4} Student Presentations}} \\
\midrule
12 & January 15 \\ % Freya 
13 & January 22 \\ % Louise
14 & January 29 \\ % Louise
15 & February 5 \\ % Louise
16 & February 12 \\ % Louise
\bottomrule
\end{tabular}
\end{center}


\end{itemize}

\end{document}
